\section{Méthode}
\subsection{Appareils de mesure}

\begin{table}[H]
    \centering
    \begin{tabular}{|l|l|l|}
        \hline
        \textbf{Type} & \textbf{Marque} & \textbf{Modèle} \\
        \hline
        Multimètre numérique & UNI-T & UT61E \\
        \hline
    \end{tabular}
    \caption{Appareils de mesure}
    \label{tab:appareils_mesure}
\end{table}


\section{Batterie}

Afin de mesurer correctement le niveau de la batterie, un pont diviseur est utilisé, comme expliqué dans le chapitre consacré à la maquette.

Il est important de s'assurer que les résistances utilisées dans le pont diviseur sont correctement reportées dans le code afin de garantir une lecture fidèle de la tension de la batterie.

\begin{table}[H]
    \centering
    \begin{tabular}{|l|c|c|}
        \hline
        \textbf{Symbole} & \textbf{Valeur théorique} & \textbf{Valeur mesurée} \\
        \hline
        R4 & 10 k$\Omega$ & 10.12 k$\Omega$ \\
        \hline
        R3 & 51 k$\Omega$ & 49.50 k$\Omega$ \\
        \hline
    \end{tabular}
    \caption{Mesure du pont diviseur}
    \label{tab:mesure_pont_diviseur}
\end{table}

 À l'aide des valeurs mesurées, le nouveau rapport du pont diviseur peut être calculé puis introduit dans le code afin d'affiner la mesure de la tension de la batterie.

Le rapport théorique du pont diviseur est donné par :

$$
\frac{R_3+R_4}{R_4}
$$

En utilisant les valeurs mesurées, on obtient :

$$
\frac{50.82 + 10.12}{10.12}=6.022
$$

Le rapport théorique valait :

$$
\frac{51+10}{10}=6.10
$$

On constate donc un écart d'environ (1\%) entre le rapport théorique et le rapport réel.

Il est alors possible de corriger les constantes utilisées dans le programme en remplaçant les valeurs nominales des résistances par les valeurs mesurées :


Le facteur de conversion utilisé pour reconstruire la tension de la batterie devient alors :

$$
\frac{50820+10120}{10120}=6.022
$$

Cette correction permet d'améliorer la précision de la mesure de la tension et de réduire l'erreur introduite par les tolérances des résistances du pont diviseur.

\section{Tensions, Courants}
Afin de mesurer la consommation énergétique de la maquette, l'idée est de réaliser une mesure du courant total dans le système.
Pour ce faire, le multimètre est placé en mode ampèremètre entre le connecteur de la batterie et celui du PCB.

\begin{table}[H]
    \centering
    \begin{tabular}{|l|c|c|}
        \hline
        \textbf{Grandeur mesurée} & \textbf{Valeur théorique} & \textbf{Valeur mesurée} \\
        \hline
        $V_{bat}$ & - & 16.48 V \\
        \hline
        $V_{IO}$ & 3.3 V & 3.28 V \\
        \hline
        $I_{tot}$ & - & 0.5 A \\
        \hline
    \end{tabular}
    \caption{Mesures électriques - Moteurs allumés}
    \label{tab:moteurs_allumes}
\end{table}

\begin{table}[H]
    \centering
    \begin{tabular}{|l|c|c|}
        \hline
        \textbf{Grandeur mesurée} & \textbf{Valeur théorique} & \textbf{Valeur mesurée} \\
        \hline
        $V_{bat}$ & - & 16.52 V \\
        \hline
        $V_{IO}$ & 3.3 V & 3.28 V \\
        \hline
        $I_{tot}$ & - & 0.5 A \\
        \hline
    \end{tabular}
    \caption{Mesures électriques - Moteurs éteints}
    \label{tab:moteurs_eteints}
\end{table}

\section{Autonomie}

\section{Autonomie}